%% APPENDIX A -- spectral response, frequency conventions, primary-beam response.
%% Owner: appendix-triage.  Carries \label{app:conventions} (cited by
%% 04_method.tex) and \label{app:pbresp}.
\section{Spectral response, frequency conventions and the primary-beam
response}
\label{app:conventions}

This appendix sets out the three things every threshold quoted in this paper
is checked from: the frame of each frequency, the instrumental response, and
the primary-beam factor that turns a measured noise into a flux at the star.


%% The single-column figure of the three-point Hanning response against
%% sub-channel phase is removed; its numbers are in the peak-channel
%% paragraph below.  QUEUE_integrate.md entry S6: make_all.sh must stop
%% running make_fig_hanning.py and figorphan.py must carry
%% hanning_response.pdf in PENDING, or the figure gate fails.

\subsection{Frequency frames, linewidth and smearing}
%% (subsection label retired: nothing cites it)

Every frequency axis here --- searched spectra, tabulated ranges and
crossing frequencies alike --- is \emph{topocentric} as delivered by the
ALMA correlator, and the search itself applies no velocity-frame shift.
Velocities enter at the vetting stage, where the line mask is evaluated in
each star's own rest frame: it carries every crossing frequency to the
barycentre at its own block's epoch and then to the stellar frame before
comparing it with a transition (Appendix~\ref{app:linemask}). The residual
error is then the catalogue radial-velocity uncertainty alone, a few
km\,s$^{-1}$ within 40\,pc; second-order relativistic terms at the
\MaskHalfKms\,km\,s$^{-1}$
scale are $\lesssim5$\,kHz at 350\,GHz, below the finest native channel
searched. The stacking is drift-aligned and inverse-variance weighted:
coherent in frequency trajectory, not in electric-field phase.

The native channel width is also the linewidth convention for every EIRP
threshold. The statistic is the single-channel flux density after drift
correction, so EIRP$_{\rm trig}$ is the minimum \emph{total} power of a
channel-confined carrier, flat in the assumed linewidth below the channel and
never a spectral power density; flatness is exact only for a centred tone,
and a boundary-straddling one divides its power between two channels, which
is the loss the injections carry and one reason EIRP values are quoted to two
significant figures. The Hz-resolution
surveys' limits are total powers in the same sense, each the best case of
its own channelisation, which makes the axes of Fig.~\ref{fig:context}(a)
commensurable. One number fixes the scale, illustratively and not
achievably: rechannelising the median Class~A window's
\SqrtScaleChanKHz\,kHz to the \SETIChanHz\,Hz of a dedicated narrowband
search would divide the threshold by $\SqrtScaleFac$, and no reprocessing
recovers resolution the correlator never recorded.

Channel widths come from the measurement sets' own CHAN\_WIDTH; MS~v2
defines no ``EFFECTIVE\_BANDWIDTH'' column, so the quoted channelisations are
separations and are not claimed to be noise bandwidths.

ALMA applies online Hanning smoothing by default \citep{ALMAHandbook}, so
the spectral response is ${\sim}2$ channels wide and neighbouring channels
are correlated. For the $0.25/0.5/0.25$ kernel the coefficients
are $\rho_1=\HanRhoOneFrac=\HanRhoOne$ and $\rho_2=\HanRhoTwo$, with
$\sigma_{\rm smoothed}/\sigma_{\rm raw}=\HanSigRatio$: the correlation is
not confined to adjacent channels, the lag-two term being a sixth. The peak
channel holds \HanBest{} of an unresolved carrier's flux at a channel centre
and \HanWorst{} on a boundary, a median \HanMed{} over sub-channel phase, so
a nominal threshold is optimistic by
$\times\HanFacBest$--$\times\HanFacWorst$, and every threshold quoted here
carries that correction window by window.

\emph{The injected carriers are deposited with the correlator's own channel
response, at a sub-channel phase drawn afresh for each tone, so the smoothing
loss sits inside the measured completeness rather than being applied to it.}
The deposition falls \InjShortPct{} per cent short of the ideal response
averaged over phase, so the completeness is conservative by that much. Where
in a channel a carrier was placed makes no difference to its recovered
signal-to-noise, because a carrier at a typical searched drift rate crosses
about \InjDriftChan{} channels during one observation and so samples the
whole response. Since the loss is carried end to end it must not be applied
to the result as well: at the ninety per cent point of the recovery curve the
peak channel holds $\InjPeakAtNinety\times$ the nominal trigger flux, which
is that factor of about two, already spent.

Two measurements fix the smoothing state. The retained per-integration
residuals of the second-stage local null have a lag-one spectral
autocorrelation of \LNLagOne{} against the Hanning prediction \HanRhoOne, so
the correlation is instrumental and those spectra are smoothed and
un-averaged; and the archive's reported effective spectral resolution is
exactly twice the channel separation in \NHanTwo{} of the \NHanKnown{}
windows, in both correlator modes; the other \NHanOther{} report a smaller
ratio, the signature of online channel averaging, where the peak-channel loss
is smaller. Every window has a reported resolution, so the blanket
factor is right for most of both classes. \emph{The correlation model is truncated at
lag two}: the baseline step, a per-integration median over up to \MedWin{}
channels, could imprint longer-range correlation that the retained spectra
carry no residuals to measure, which would make the independent-trials count
of Appendix~\ref{app:falsealarm} an over-estimate.

\emph{Transmitters wider than one channel.} The baseline filter
(\S\ref{sec:statistic}) is flat for features narrower than roughly half its
effective width $W$ and falls to zero above $W$, so the survey's sensitivity
is to channel-confined morphologies and $W$ is the linewidth ceiling:
\MedWinFineLo--\MedWinFineHi{} channels over the fine class and
\MedWinCoarseLo--\MedWinCoarseHi{} over the coarse one, that is
${\approx}\MedWinCeilFineMHz$\,MHz at \ChanFineKHz\,kHz and
${\approx}\MedWinCeilCoarseGHz$\,GHz at \ChanCoarseMHz\,MHz. Within that
ceiling a top-hat transmitter spanning $N_{\rm ch}$ native channels has
per-channel flux the total divided by $N_{\rm ch}$ at unchanged noise, so
its threshold is $N_{\rm ch}$ times the single-channel one.

\emph{Intra-integration smearing, which the injections do not model.} Each
injected tone is deposited at one frequency per integration, so the campaign
carries the channel response and the drift between integrations but not the
sweep within one. A carrier moving $\Delta$ channels during an integration
keeps $\eta_{\rm smear}=\mathrm{sinc}(\pi\Delta/2)$ of its peak amplitude, and
every limit in this paper is divided by it. Since
$\mathrm{EIRP}_{90}$ is a ninety per cent point over carriers drawn from each
window's own drift grid, the $\Delta$ that sets it is the one at the ninetieth
percentile of that draw, \SmcQninePct{} per cent of the searched ceiling over
the campaign's \SmcNTone{} tones. The integration length is recorded for
\SmcNMeas{} of the \NWindows{} windows, \CsNDumpWin{} of them in the delivered
metadata and \SmNRecov{} more in the exposure records the search kept; for the
other \SmcNBound{} the factor is \emph{bounded} at \SmcDtCeil\,s, the longest
integration delivered anywhere, where none falls below \SmcBoundEta.
Nothing moves in the median window,
and \SmcNAbove{} windows on \SmcNAboveStar{} stars rise by more than one per cent,
up to $\times\SmcFacWorst$: \SmcAboveStars. Re-scoring each of the
\SmcValNWin{} injected windows' own recovery ladders with every tone de-rated
by its own $\eta_{\rm smear}$ gives a penalty never larger than the one
applied and smaller by at most \SmcValWorstPct{} per cent, so the correction
errs the safe way.

\subsection{The primary-beam response, and the two columns that let a reader
recompute it}
%% (subsection label retired: nothing cites it)

The sensitivity quoted for a window is $S_{\min}=5\sigma/A$, where $A$ is the
primary-beam response at the star's offset from the phase centre: the noise
$\sigma$ is measured in apparent flux, the limit is quoted in true flux, and
the two differ by $A$. That offset is not a nuisance. Archival fields are
pointed at a catalogue position, often decades old, while these are
high-proper-motion stars, so the star can sit well off axis: over the
\PbNWin{} windows the offset has a median of \PbOffMed$''$, a 99th
percentile of \PbOffPNN$''$ and a maximum of \PbOffMax$''$, and $A$ runs from
\PbAMed{} in the median down to \PbAMin, below $0.99$ on \PbNLowNN{} windows
and below $0.9$ on \PbNLowN. \emph{None} falls below one half, and every one
of the low-response windows is a physically explicable J2000 pointing on a
fast-moving star rather than an error. A window is kept only where the
response is at least \PbFloorResp, a correction of at most
$\times\PbFloorCorr$, since at low response a small error in the beam model
becomes a large error in inferred flux. The offsets are bimodal and the floor
sits in the gap: the last retained window has response \PbGapKeptResp{} at
$\PbGapKeptU\,\theta_{\rm PB}$, the first window beyond it \PbGapHeldResp{}
at $\PbGapHeldU\,\theta_{\rm PB}$, and nothing lies between, so any floor
between those two responses retains the same \NWindows{} windows and the
value chosen does not matter.
Beyond it lie the \EpsEriNWin{} $\epsilon$\,Eridani Band~6 windows, which
were withheld as a band and not window by window. \EpsEriNOffAxis{} of them,
in \LqEpsOffEb{} blocks, need a correction of up to
$\times\EpsEriCorrMax$ --- past the first sidelobe transition, where the
Gaussian form fails --- and so fail the floor outright. The other
\EpsEriNOnAxis, in block \texttt{\BkEpsOnBlock}, place the star at the
pointing centre and need no correction at all, so the floor is not what
removed them; and at \EpsEriDistPc\,pc the star is too near to be left
without a limit, so theirs are given here. They carry
\BkEpsOnHours\,window-hours and no threshold crossing, the largest statistic
being \BkEpsOnTMax{} where the trigger is five. The one fine-channel window,
at \BkEpsOnChanFineKHz\,kHz, reaches $\BkEpsOnEirpFine$\,W and the
\BkEpsOnNCoarse{} coarse ones at \BkEpsOnChanCoarseMHz\,MHz
$\BkEpsOnEirpCoarseLo$--$\BkEpsOnEirpCoarseHi$\,W; none enters a census
statistic. The flux-scale comparison of
\S\ref{sec:method} is a separate continuum measurement on the same star: a
block mean over the \PxEpsN{} compact-array executions that place it at the
pointing centre, \PxEpsHours\,h in Band~\PxEpsBand, where no beam correction
enters either.

\emph{The correction is applied exactly once.} The offset is computed for
every measurement set from its own \texttt{FIELD::PHASE\_DIR} and the
propagated stellar position, so no window carries a zero-by-assumption
offset, and it enters in one place only: $S_{\min}=5\sigma/A$ closes on all
\PbSminClosed{} catalogued rows to \PbSminWorstRes{} in relative terms. The
count is in two parts, because on \PbSminTaut{} rows whose products are not
retained $A$ is back-solved from $5\sigma/S_{\min}$ and flagged as such, so
there the closure is arithmetic; the audit is a test on the other
\PbSminIndep{} rows, where $A$ comes from the search's own stored geometry and
closes to \PbSminIndepRes. Nothing else needs correcting: the statistic of
Eq.~\ref{eq:tstar}, the control ranks and the attribution are \emph{ratios}
formed inside one window, in which $A$ divides out identically.

\emph{The response applied is a Gaussian} of full width
$\RingPbCoefCode\,\lambda/D$ --- an Airy first-null coefficient --- at each
window's \emph{own} dish diameter, reproduced over all \NWindows{}
rows to $\SensPbResid$ against $\SensPbResidTwelve$ for the same Gaussian
forced to $D=12$\,m. The distinction is not cosmetic: the \PbNAcaRow{} rows
taken on the \PbAcaRatio$\times$ smaller compact-array dishes have a
correspondingly wider beam, and their response cannot be obtained from a
12\,m beamwidth at all.

\emph{The beamwidth convention is a bounded systematic, and both
conventions were computed.} ALMA's
primary-beam full width is $\RingPbCoef\,\lambda/D$, narrower than the
first-null coefficient, so it attenuates more and makes every limit
\emph{shallower}: by nothing at all in the median window, by more than one
per cent in \SensPbNOnePct{} of \NWindows{} and by more than ten in
\SensPbNTenPct, the largest being $+\SensPbShiftMaxPct$ per cent --- those
few all \SensPbWorstStar{} at \SensPbWorstOff$''$ off axis. The lowest
response in the sample moves \SensPbAMin{} to \SensPbAMinAlma, still clear
of the \SensPbFloor{} retention floor but by a margin worth saying.
No disposition changes with the convention. The
remaining model uncertainty is also bounded by measurement: CASA's
blocked-Airy response is better physics than either Gaussian, the largest
ratio to it over the whole sample is \PbModelMax, and the counts of systems
reaching $10^{15}$\,W are identical under any of the three.

\emph{The star's apparent direction includes the annual parallax, and the
amplitude a compact source retains when that is left out of the assumed phase
centre is measured rather than modelled.} The displacement decides whether the
loss can be treated as a scalar at all: a small fraction of a synthesised beam
attenuates a compact source and does nothing else, while a displacement
approaching a beam puts the phase wrong on the longest baselines. Re-extracting
the \SensNReextWin{} windows of \SensNReextBlock{} blocks furthest out at the
full apparent position --- proper motion, parallax and radial velocity where
catalogued --- leaves each window's noise unchanged to
\SensReextNoiseMaxPct{} per cent and alters no disposition. It cannot test
the amplitude factor, because that factor describes a source and those
windows hold none: the statistic is a maximum over a plane of drift
trajectories, so moving the extraction point by two thirds of a beam draws a
fresh noise realisation, and the ratio scatters from
\SensReTstarLo{} to \SensReTstarHi{} about a median of \SensReTstarMed{}
instead of approaching the \SensPxLargePred{} a source would give. A
carrier injected at the apparent position and recovered from the \emph{same}
visibilities at both positions does return the normalised dirty beam at the
omitted displacement, measured on that block's own baselines with no model in
it: \SensPxSmallMeas{} against a predicted \SensPxSmallPred{} at
\SensPxSmallBeams{} of a beam, and \SensPxLargeMeas{} against
\SensPxLargePred{} at \SensPxLargeBeams. That it exceeds the model at the
larger displacement is the search's doing, since maximising over a plane of
trajectories lets a partly decohered carrier still find a cell, so a limit
corrected by the model is shallower than it needs to be.

\emph{The injection campaign is blind to $A$ by construction, and the bias
that follows is exactly zero.} The injector deposits an unattenuated amplitude
at the star's own direction cosines and refers the ladder to the
\emph{apparent} $5\sigma$ on all \PbInjNUnit{} scored units, so an injected
apparent amplitude recovered nine times in ten implies a true flux larger by
$1/A$ while $P_{\rm trig}$ is the power of a true $5\sigma/A$ and $A$ cancels
identically; \PbInjNDistinct{} of the units have apparent and corrected
$S_{\min}$ differing by more than a part in a thousand, so the cancellation
is tested rather than vacuous. Being self-consistent in apparent flux, the
campaign cannot falsify the beam model; and it reaches only
$A\ge\PbInjAMin$, whereas \PbNBeyond{} windows on \PbNBeyondStar{} stars lie
further off axis, down to $A=\PbBeyondAMin$, so a limit quoted for one of
those rests on an extrapolated curve; their \PbNBeyondCross{} crossings are
ratios and unaffected.

\emph{A few rows carry a catalogued noise that no retained product
reproduces}, and they are declared rather than reconciled: on \PbGapN{}
windows of \PbGapStar{} in block \texttt{\PbGapEb} the catalogued $\sigma$ is
\PbGapPctLo--\PbGapPctHi{} per cent \emph{higher} than the surviving
product's, so the limits there are weaker and not stronger. With the
\PbGapNoProd{} windows whose products are not retained the declared
provenance gap is \PbGapDeclared{} windows; the attenuation columns are
unaffected.

%% Carried forward: the UV Ceti worked example (flux -> EIRP) is out of this
%% appendix because two of its four numbers had no macro behind them.  If a
%% worked example is wanted back it needs three: the flux density, the implied
%% EIRP and the distance used.  It is the clearest single verification aid the
%% paper could carry.  \label{app:repro} and the transmitter-power benchmark
%% are both resolved: nothing cites the label any more, and exactly one
%% benchmark survives, in Sec. 6.
%% tab:blockledger lives here rather than in Sec. 3: it is a full-width
%% table* worth 0.53 pp of the main text, and Sec. 3 carries its headline
%% rows in prose and cites it.
\begin{table*}
\centering
\scriptsize
\caption{\textbf{The survey ledger, part one: execution blocks and windows.}
Each row carries its own number, and a group sums to the row whose number
heads it --- rows 3.1.1 to 3.1.4 to row 3.1, rows 3.1 to 3.5 to row 3. The
primary census is the survey: unless a statement names one of the other sets,
it is computed over that row. Hours are window-hours: on-source time summed
over spectral windows, so the classes partition the hours as they partition
the windows. The telescope on-source time behind the primary census is
\StTelHours\,h, of which \StTelHoursFine\,h carries at least one Class~A
window and \StTelHoursCoarse\,h at least one Class~B window --- figures that
overlap and must not be added. Confirmed planet counts here and in the
target list come from the \SfHostCatalogue{} \SfHostTable{} table
\citep{NEA2013}, accessed \SfHostDate. Part two, Table~\ref{tab:counts},
carries the events and the confirmation extents.}
\label{tab:blockledger}
%% GENERATED by blockledger_v411.py -- do not hand-edit.
\begin{tabular}{@{}r@{~}p{0.265\textwidth}rrrrrp{0.265\textwidth}@{}}
\hline
 & Set of execution blocks & Blocks & Windows & Stars & Systems & Hours & Results it enters \\
\hline
1 & Archival scope: blocks the query exposes & 656 & -- & -- & -- & -- & rows 1.1 and 1.2 \\
1.1 & \hspace{0.8em}\ldots{} processed with the frozen pipeline & 479 & -- & -- & -- & -- & panel 2 \\
1.2 & \hspace{0.8em}\ldots{} not processed at the freeze & 177 & -- & -- & -- & -- & panel 3 \\
\hline
2 & Processed with the frozen pipeline$^{c}$ & 484 & -- & -- & -- & -- & rows 2.1 to 2.4 \\
2.1 & \hspace{0.8em}\textbf{Primary census} & 404 & 1651 & 89 & 82 & 1094 & every limit, the crossing list, every chance expectation \\
2.1.1 & \hspace{1.6em}Class~A, fine channels$^{a}$ & -- & 402 & 65 & 60 & 227 & the drifting-carrier experiment \\
2.1.2 & \hspace{1.6em}Class~B, coarse channels$^{a}$ & -- & 1249 & 88 & 81 & 867 & the spectral-excess experiment \\
2.2 & \hspace{0.8em}Reserved hold-out & 77 & 315 & 36 & -- & 214 & the rank displacement, the false-alarm tail factor, the radial control profile \\
2.3 & \hspace{0.8em}Withheld: primary-beam form invalid$^{d}$ & 2 & 8 & -- & -- & -- & nothing; withheld \\
2.4 & \hspace{0.8em}$\epsilon$\,Eri at the pointing centre$^{e}$ & 1 & 4 & 1 & 1 & 0.1 & reported limits only; no census statistic \\
\hline
3 & Not processed at the freeze: what became of each & 177 & -- & -- & -- & -- & rows 3.1 to 3.5, the five fates \\
3.1 & \hspace{0.8em}Calibrated and searched & 144 & -- & -- & -- & -- & rows 3.1.1 to 3.1.4 \\
3.1.1 & \hspace{1.6em}\textbf{Archival extension}, reported & 115 & 463 & 16 & 16 & 269 & recurrence, and the null on data the frozen pipeline had not seen \\
3.1.2 & \hspace{1.6em}Searched after the extension was frozen & 8 & 32 & -- & -- & -- & nothing; later than the reported extension \\
3.1.3 & \hspace{1.6em}Withheld: primary-beam form invalid$^{d}$ & 8 & 32 & -- & -- & -- & nothing; withheld \\
3.1.4 & \hspace{1.6em}No window reached the export & 13 & 0 & -- & -- & -- & nothing \\
3.2 & \hspace{0.8em}No calibration in the ALMA delivery & 18 & -- & -- & -- & -- & nothing; the archival gap \\
3.3 & \hspace{0.8em}Calibration failed & 11 & -- & -- & -- & -- & nothing; the archival gap \\
3.4 & \hspace{0.8em}Not attempted & 3 & -- & -- & -- & -- & nothing; the archival gap \\
3.5 & \hspace{0.8em}Working disk insufficient & 1 & -- & -- & -- & -- & nothing; the archival gap \\
\hline
4 & Beyond 40\,pc: out-of-sample control & 47 & 198 & 22 & -- & 88 & the external null; no science count \\
5 & External calibration sample$^{b}$ & 326 & 1322 & 56 & -- & -- & the rank calibration and the positive control; no science count \\
\hline
\end{tabular}

{\footnotesize $^{a}$The two classes partition the census windows, not its blocks: 252 blocks carry a Class~A window and 389 a Class~B one, and most carry both, so the hours of the two rows partition the windows and not the observing time.\\ $^{b}$Not an out-of-sample set and not an addition to this ledger: every one of its 326 blocks is counted in a row above, 263 in the primary census and 63 in the hold-out, and it is listed apart only to keep it from being counted twice. One further block of the harvest it came from lies outside the 40\,pc work list and is excluded here.\\ $^{c}$Panels 1 and 2 overlap rather than add: 479 of these 484 blocks are row 1.1, and the other 5 carry no member observing unit in the archive metadata, so the scope query does not expose them; 3 of the 5 are in the primary census.\\ $^{d}$Every one is an $\epsilon$\,Eri Band~6 block, and the band was withheld whole: in most of them the star sits far enough off axis (response 0.28--0.35, below the 0.5 floor) that the Gaussian primary-beam form is not valid.\\ $^{e}$The windows of the same band that place the star at the pointing centre, at a response of 1.00, and so pass the floor; their limits are in Appendix~\ref{app:conventions} and they enter no census statistic. Row 3.1.1 holds a further 16 such windows, which enter nothing either.}

\end{table*}

\begin{table}
\centering
\scriptsize
\caption{\textbf{The survey ledger, part two: events, dispositions and
confirmation.} Every headline count in this paper is read from here. Three
counts are \MtCrossA{} over three different sets --- the Class~A crossings of
the adopted ledger, those of the extraction catalogue, and the
quality-passing crossings. The first two hold different rows, since re-extraction moved four
crossings across the trigger each way; the last two hold the same rows,
because the one block failing the quality criterion is also the survey's only
block with a coarse-channel crossing. Row 4.2 is what this
paper means by confirmation and the only such count it quotes: a repeat epoch
at a frequency the first did not cover confirms nothing. Row 2.2.1 was also
decomposed system by system, each expectation formed from that system's own
control positions: \MtNSysLowP{} of the \MtNSysChance{} systems fall below
one chance in twenty, against \MtNSysLowPExp{} expected, and \MtTopStar,
which holds \MtTopStarNCross{} of the \EvNCrossRaw{} crossings,
\MtTopStarNUnattr{} of them unattributed, gives \MtTopSysObs{} in census
windows against \MtTopSysExp{} expected over its \MtTopSysNWin{} Class~A
windows. No system
departs from its own controls by more than chance allows.}
\label{tab:counts}
%% GENERATED by counts_r15.py -- do not hand-edit.
\begin{tabular}{@{}r@{~}p{0.405\columnwidth}r@{\,\,}p{0.325\columnwidth}@{}}
\hline
 & Count & Value & What it counts \\
\hline
1 & Threshold crossings at $T_\star\ge5$ & \EvNCrossRaw & census and archival tail \\
1.1 & Class~A, fine channels & \MtCrossA & the carrier experiment \\
1.1.1 & \ldots{} in the primary census & \MtCrossACen & catalogued window by window \\
1.1.2 & \ldots{} in archival-tail blocks & \CeNTailCross & both toward \CeTailStar; in no expectation \\
1.2 & Class~B, coarse channels & \LdgNCrossB & all in the one block failing \S\ref{sec:bdblock} \\
\hline
2 & Quality-passing crossings & \EvNCrossQp & the same rows as 1.1 \\
2.1 & \ldots{} attributed to a transition & \EvNAttrQp & inside the mask in the star's own frame \\
2.2 & \ldots{} unattributed & \EvNUnattrQp & rows 2.2.1 and 2.2.2 \\
2.2.1 & \ldots{} in the primary census & \CeObs & against \CeExp{} expected by chance, measured at the \CeNWin{} windows' own control positions (\CeLo--\CeHi{} allowed) \\
2.2.2 & \ldots{} in archival-tail blocks & \CeNTailCross & outside that comparison \\
2.2.3 & \ldots{} above all \NCtrl{} controls & \EvNRankQp & the two events of Fig.~\ref{fig:event} \\
2.2.4 & \ldots{} recovered in a later block & \NUnattNoRepeat & no confirmed signal; \MkAttrRecur{} attributed ones do \\
\hline
3 & Class~A systems & \NSysClassA & of \NSystems{} surveyed \\
3.1 & \ldots{} seen again after \EpSep{} day & \ScSysIndep & at any frequency \\
3.2 & \ldots{} at the same stellar-frame frequency & \RcConfSysDay & \textbf{where a signal could have been confirmed}; \RcConfSysYr{} of them over a year later \\
3.3 & \ldots{} with no such second block & \RcNoRepSys & confirmation impossible; \RcNoRepSysAll{} of all \NSystems \\
\hline
4 & Class~A union bandwidth & \CvUnionA\,GHz & \DomAIslands{} disjoint islands \\
4.1 & \ldots{} removed by the line mask & \MkMaskAPct\,\% & \MkMaskAGHz\,GHz: the figure the text quotes, against \FrMaskPct{} per cent of the \DnuAB\,GHz archival union \\
\hline
\end{tabular}

\end{table}


%% ★★ 2026-10-08: THE ALMA PROJECT CODES, RESTORED AS A TABLE.
%% They were taken out of the Acknowledgements as padding and then lived "in
%% the data release"; there is no data release, so they were nowhere.  ALMA's
%% data-use policy expects the codes to be traceable from the publication, so
%% this is not optional -- but an appendix table satisfies the policy without
%% putting \NProjCodes{} codes back into the front matter, which is what was
%% objected to.  The fragment is generated by projcodes_v400.py and is the
%% same list that resolves every block in Table~\ref{tab:blockledger}.
\begin{table}
\centering
\scriptsize
\caption{\textbf{The \NProjCodes{} ALMA proposals the searched execution
blocks belong to}, in the standard
\mbox{ADS/JAO.ALMA\#\textit{code}} form. Every block of the primary census
belongs to one of these, and the blocks themselves are retrieved from the
ALMA Science Archive by the execution-block identifiers
Tables~\ref{tab:ledger} and~\ref{tab:ledgercont}
give.\ProjCodeUnresClause}
\label{tab:projcodes}
%% GENERATED by projcodes_v400.py -- do not hand-edit.
\mbox{2012.1.00142.S}, \mbox{2012.1.00198.S}, \mbox{2012.1.00268.S}, \mbox{2012.1.00385.S}, \mbox{2012.1.00437.S}, \mbox{2012.1.00993.S}, \mbox{2013.1.00359.S}, \mbox{2013.1.00588.S}, \mbox{2013.1.01147.S}, \mbox{2013.1.01214.S}, \mbox{2015.1.00307.S}, \mbox{2015.1.00633.S}, \mbox{2015.1.00783.S}, \mbox{2015.1.01015.S}, \mbox{2015.1.01260.S}, \mbox{2016.1.00104.S}, \mbox{2016.1.00358.S}, \mbox{2016.1.00515.S}, \mbox{2016.1.00637.S}, \mbox{2016.1.00803.S}, \mbox{2016.1.00817.S}, \mbox{2016.1.00878.S}, \mbox{2016.1.00880.S}, \mbox{2016.2.00015.S}, \mbox{2016.2.00200.S}, \mbox{2016.A.00013.S}, \mbox{2017.1.00167.S}, \mbox{2017.1.00200.S}, \mbox{2017.1.00215.S}, \mbox{2017.1.00561.S}, \mbox{2017.1.00698.S}, \mbox{2017.1.00704.S}, \mbox{2017.1.00786.S}, \mbox{2017.1.00986.S}, \mbox{2017.1.01583.S}, \mbox{2017.1.01644.S}, \mbox{2018.1.00045.S}, \mbox{2018.1.00072.S}, \mbox{2018.1.00126.S}, \mbox{2018.1.00461.S}, \mbox{2018.1.00470.S}, \mbox{2018.1.00947.S}, \mbox{2018.1.01149.S}, \mbox{2018.1.01545.S}, \mbox{2018.1.01833.S}, \mbox{2019.1.00189.S}, \mbox{2019.1.01443.T}, \mbox{2019.1.01517.S}, \mbox{2019.1.01603.S}, \mbox{2019.1.01681.S}, \mbox{2019.2.00208.S}, \mbox{2021.1.00485.S}, \mbox{2021.1.01209.S}, \mbox{2021.2.00097.S}, \mbox{2022.1.00134.S}, \mbox{2022.1.00338.L}, \mbox{2022.1.00793.S}, \mbox{2022.1.01163.S}, \mbox{2023.1.00390.S}, \mbox{2023.1.00487.S}, \mbox{2024.1.00165.S}, \mbox{2024.1.00595.S}, \mbox{2024.1.00722.S}, \mbox{2024.1.01095.S}, \mbox{2024.A.00040.S}

\end{table}
